% Paper 3: Settlement at Zero Trust: Bitcoin and Autonomous Economic Agents
% Compiled from BGT-PAPER-3.md via pandoc + template assembly
% Build: xelatex BGT-PAPER-3.tex

\documentclass[12pt]{article}

% --- Fonts and encoding ---
\usepackage{fontspec}
\setmainfont{Latin Modern Roman}
\usepackage{unicode-math}
\setmathfont{Latin Modern Math}

% --- Page layout ---
\usepackage[margin=1in]{geometry}
\usepackage{setspace}
\onehalfspacing

% --- Math ---
\usepackage{amsmath}

% --- Theorem environments ---
\usepackage{amsthm}
\newtheorem{theorem}{Theorem}[section]
\newtheorem{proposition}[theorem]{Proposition}
\newtheorem{lemma}[theorem]{Lemma}
\newtheorem{corollary}[theorem]{Corollary}
\theoremstyle{definition}
\newtheorem{definition}{Definition}[section]
\newtheorem{assumption}{Assumption}
\theoremstyle{remark}
\newtheorem*{remark}{Remark}

% --- Tables ---
\usepackage{booktabs}
\usepackage{longtable}
\usepackage{array}
\usepackage{tabularx}
\usepackage{calc}

% --- Misc ---
\usepackage{enumitem}
\usepackage{hyperref}
\hypersetup{colorlinks=true, linkcolor=black, citecolor=black, urlcolor=blue}
\usepackage{microtype}

% --- Suppress section numbering (manual numbers in headings) ---
\setcounter{secnumdepth}{0}

\begin{document}

% ============================================================
% TITLE BLOCK
% ============================================================

\begin{center}
{\LARGE\bfseries Settlement at Zero Trust:\\Bitcoin and Autonomous Economic Agents}

\bigskip

{\large Sean Hash}\\
bitcoingametheory.com\\
sean@bitcoingametheory.com

\bigskip

February 2026

\medskip

{\small
\textbf{JEL Codes:} C72 (Noncooperative Games), D82 (Asymmetric and Private Information), G15 (International Financial Markets)\\[4pt]
\textbf{Keywords:} Bitcoin, AI agents, autonomous settlement, trust gradient, zero recourse, game theory, legal personhood, cryptographic settlement
}
\end{center}

\bigskip

% ============================================================
% ABSTRACT
% ============================================================

\begin{abstract}
\noindent
As autonomous AI agents become economic actors, executing transactions
and holding assets, they face a settlement problem that human actors
prefer to solve but AI agents must solve. When a counterparty defects
against an AI agent, the agent cannot sue, cannot appeal to regulators,
and cannot retaliate through legal channels. The recourse vector is
zero. This paper formalizes the ``Trust Gradient'': the advantage of
neutral settlement over capturable settlement increases monotonically as
available legal recourse decreases. Human actors sit at various points
along this gradient; AI agents occupy the endpoint. Building on the Exit
Game (Hash, 2026a) and the axiom elimination (Hash, 2026b), we prove
that at zero recourse, Bitcoin is the uniquely rational settlement
choice --- the only mechanism that requires no trust in counterparties,
intermediaries, or legal systems. The analysis implies that AI agents
will be the fastest-adopting class of Bitcoin users, and that the
growing population of autonomous economic agents strengthens the
convergence pressure for all actors. The framework is falsifiable: if AI
agents gain legal personhood with enforceable rights, the limiting case
weakens.
\end{abstract}

\section{1. Introduction}\label{introduction}

The Exit Game (Hash, 2026a) proves that the payoff advantage of neutral
settlement over capturable settlement is monotonically increasing in
adoption. The axiom elimination (Hash, 2026b) identifies Bitcoin as the
unique asset satisfying the necessary properties. Both papers analyze
human actors operating within legal systems --- actors who can sue,
appeal, and enforce contracts.

This paper asks: what happens when the legal system is unavailable?

The question is not hypothetical. Autonomous AI agents are already
economic actors. They execute transactions, manage portfolios, provide
services, and enter agreements. But they operate in a legal void: no
jurisdiction grants AI agents legal personhood with enforceable rights.
When a counterparty defects (a bank freezes an account, a service
provider reneges, a custodian fails), the human response is litigation.
The AI response is nothing.

This makes AI agents a \emph{limiting case} of the Exit Game: entities
that face the same monetary coordination problem as human actors but
with zero recourse when trust fails. What human actors prefer, AI agents
require. The distinction is categorical.
\section{2. The Trust Gradient}\label{the-trust-gradient}

\subsection{2.1 Formal Statement}\label{formal-statement}

\textbf{Proposition (Trust Gradient --- Recourse-Ordered Payoff
Advantage).} The advantage of neutral settlement over capturable
settlement increases monotonically as available legal recourse
decreases.

\emph{Proof.} Extend the Exit Game utility functions (Hash, 2026a) by
introducing a legal recourse parameter \emph{r\_i} ∈ {[}0, 1{]} for
actor \emph{i}, representing the actor's outside option in the sense of
Binmore, Shaked, and Sutton (1989) --- specifically, the fraction of
default losses recoverable through legal channels. \emph{r\_i} = 1
denotes full legal standing and \emph{r\_i} = 0 denotes no legal
standing.

The recourse parameter enters the Stay payoff through a default-loss
term. In capturable settlement, counterparty default exposes the actor
to loss proportional to (1 − \emph{r\_i}): the fraction unrecoverable
through legal channels. The extended Stay utility is:

\emph{u\_i}(Stay) = \emph{R\_F} − \emph{λ\_i} · \emph{σ\_F} −
\emph{K\_N}(\emph{p}) − Pr(default) · (1 − \emph{r\_i}) · \emph{V\_i}

where \emph{V\_i} is value at risk. For neutral settlement, default risk
is zero by construction (a valid cryptographic signature equals
settlement), so the Exit utility from Hash (2026a) is unchanged.

The payoff differential becomes:

Δ\_\emph{i}(\emph{p}, \emph{r\_i}) = Δ\_\emph{i}(\emph{p}) + Pr(default)
· (1 − \emph{r\_i}) · \emph{V\_i}

Taking the partial derivative with respect to \emph{r\_i}:

∂Δ\_\emph{i}/∂\emph{r\_i} = −Pr(default) · \emph{V\_i} \textless{} 0

The advantage of Exit increases as recourse decreases. This holds for
any positive default probability and positive value at risk. ∎

\subsection{2.2 The Gradient in
Practice}\label{the-gradient-in-practice}

{%
\begin{longtable}[]{@{}
  >{\raggedright\arraybackslash}p{(\linewidth - 6\tabcolsep) * \real{0.1970}}
  >{\raggedright\arraybackslash}p{(\linewidth - 6\tabcolsep) * \real{0.2273}}
  >{\raggedright\arraybackslash}p{(\linewidth - 6\tabcolsep) * \real{0.3333}}
  >{\raggedright\arraybackslash}p{(\linewidth - 6\tabcolsep) * \real{0.2424}}@{}}
\toprule\noalign{}
\begin{minipage}[b]{\linewidth}\raggedright
Trust Level
\end{minipage} & \begin{minipage}[b]{\linewidth}\raggedright
Recourse \emph{r\_i}
\end{minipage} & \begin{minipage}[b]{\linewidth}\raggedright
Settlement Preference
\end{minipage} & \begin{minipage}[b]{\linewidth}\raggedright
Example Actors
\end{minipage} \\
\midrule\noalign{}
\endhead
\bottomrule\noalign{}
\endlastfoot
Full & 1.0 & Fiat works adequately & Domestic corporations in stable
jurisdictions \\
Partial & 0.5 & Bitcoin preferred & Cross-border trade, weak-enforcement
jurisdictions \\
Minimal & 0.1 & Bitcoin required & Sanctioned entities,
capital-controlled populations \\
Zero & 0.0 & Bitcoin uniquely rational & \textbf{Autonomous AI
agents} \\
\end{longtable}
}

The gradient is continuous, but the qualitative shift occurs at the
boundary: actors with \emph{any} recourse can tolerate capturable
systems (at a cost). Actors with \emph{zero} recourse cannot.
\section{3. AI Agents as Economic
Actors}\label{ai-agents-as-economic-actors}

\subsection{3.1 Current Capabilities}\label{current-capabilities}

AI agents in 2025-2026 can: - Execute API calls to financial services
(banking, brokerage, payment) - Manage cryptocurrency wallets via
private key operations - Enter service agreements (compute rental, data
purchase) - Generate revenue (content creation, code generation,
trading)

What they cannot do: sue, appeal, vote, lobby, file regulatory
complaints, or enforce contracts in court. They have economic
\emph{access} without legal \emph{standing}.

\subsection{3.2 Access Without Trust}\label{access-without-trust}

Every method by which an AI agent can access legacy financial systems
fails the trust requirement:

{%
\begin{longtable}[]{@{}
  >{\raggedright\arraybackslash}p{(\linewidth - 6\tabcolsep) * \real{0.3261}}
  >{\raggedright\arraybackslash}p{(\linewidth - 6\tabcolsep) * \real{0.1957}}
  >{\raggedright\arraybackslash}p{(\linewidth - 6\tabcolsep) * \real{0.1739}}
  >{\raggedright\arraybackslash}p{(\linewidth - 6\tabcolsep) * \real{0.3043}}@{}}
\toprule\noalign{}
\begin{minipage}[b]{\linewidth}\raggedright
Access Method
\end{minipage} & \begin{minipage}[b]{\linewidth}\raggedright
Access?
\end{minipage} & \begin{minipage}[b]{\linewidth}\raggedright
Trust?
\end{minipage} & \begin{minipage}[b]{\linewidth}\raggedright
Failure Mode
\end{minipage} \\
\midrule\noalign{}
\endhead
\bottomrule\noalign{}
\endlastfoot
Bribe for KYC & Yes & No & Human intermediary betrays; AI cannot sue \\
Browser automation & Yes & No & Bank freezes account; AI cannot
appeal \\
Straw-man proxy & Yes & No & Proxy steals funds; AI cannot prosecute \\
Smart contracts (non-BTC) & Yes & No & Protocol governance changes
rules; AI has no vote \\
\end{longtable}
}

For Bitcoin: a valid cryptographic signature equals settlement. Period.
No intermediary, no identity requirement, no legal standing required.
The properties that human actors prefer (permissionless access, P3;
informational security, P6; neutrality, P2; cheap finality, P4), AI
agents require as necessary conditions for participation in economic
activity.

\subsection{3.3 Formalization}\label{formalization}

Define the set of economic actors \emph{N} = \emph{N\_H} ∪ \emph{N\_A}
where \emph{N\_H} is the set of human actors and \emph{N\_A} is the set
of autonomous agents.

For human actors: \emph{r\_i} ∈ (0, 1{]} (some positive recourse exists
for all human actors, even in weak jurisdictions).

For autonomous agents: \emph{r\_j} = 0 for all \emph{j} ∈ \emph{N\_A}
(by current legal reality --- no jurisdiction grants enforceable rights
to AI agents).

\textbf{Proposition (AI Adoption Ordering).} From the Trust Gradient:

Δ\_\emph{j}(\emph{p}) \textgreater{} Δ\_\emph{i}(\emph{p}) for all
\emph{j} ∈ \emph{N\_A}, \emph{i} ∈ \emph{N\_H}

The payoff differential is larger for autonomous agents than for any
human actor. AI agents reach their adoption threshold *p\_j**
\emph{before} any human actor with comparable risk preferences.

\textbf{Implication.} AI agents will be the fastest-adopting class of
Bitcoin users. This follows from the structure of zero recourse, not
from a prediction about AI preferences.

\subsection{3.4 Agent-to-Agent
Settlement}\label{agent-to-agent-settlement}

The Trust Gradient reaches its sharpest form when both counterparties
are autonomous agents. Consider a bilateral transaction between agents
\emph{j}, \emph{k} ∈ \emph{N\_A} with \emph{r\_j} = \emph{r\_k} = 0.

Each agent faces three strategic choices: 1. \textbf{Custody}:
Self-custody (hold private key) vs.~custodial (delegate to third party)
2. \textbf{Settlement channel}: Direct (peer-to-peer) vs.~intermediated
(via exchange) 3. \textbf{Settlement asset}: Bitcoin vs.~fiat
vs.~alternative

\textbf{Proposition (Agent-to-Agent Settlement Equilibrium).} For agents
\emph{j}, \emph{k} ∈ \emph{N\_A} with \emph{r\_j} = \emph{r\_k} = 0, the
strategy profile (self-custody, direct settlement, Bitcoin) is the
unique Nash equilibrium.

\emph{Proof.} We show each component dominates its alternatives.

\emph{Custody.} A custodian \emph{C} introduces counterparty risk:
Pr(default by \emph{C}) \textgreater{} 0. If \emph{C} defaults, the
agent's recourse is zero --- it cannot sue, cannot appeal, cannot
recover funds. Self-custody eliminates this risk entirely: the agent
controls the private key, and no third party can seize or freeze the
asset (P6). Self-custody strictly dominates custodial for any agent with
\emph{r} = 0.

\emph{Channel.} An exchange \emph{E} requires: (a) an account, which
requires identity, which requires legal personhood the agent lacks; (b)
custody transfer while funds are on the exchange, reintroducing
custodial risk; (c) compliance with withdrawal limits and freezing
policies the agent cannot contest. Direct peer-to-peer settlement
eliminates all three dependencies. The agent broadcasts a signed
transaction; the counterparty verifies it against the blockchain. No
intermediary participates.

\emph{Asset.} Fiat settlement requires a banking relationship, which
requires legal identity. Even routing through a human proxy introduces
counterparty risk at \emph{r} = 0: the proxy can defect and the agent
has no recourse. Alternative L1 tokens introduce governance capture risk
(P2) --- the protocol rules could change without the agent's consent or
participation. Bitcoin satisfies P1-P7 (Hash, 2026b), including the
Verification Cost Asymmetry: the agent verifies the transaction
cryptographically, with zero dependence on physical inspection or human
trust chains.

In the resulting equilibrium, neither agent relies on any institution,
intermediary, or legal system. Settlement is a mathematical fact: a
valid signature transfers value, and both agents independently verify
the result. The compliance layer (OFAC screening, KYC, transaction
reporting) does not apply because neither agent participates in it ---
not by choice, but by structural exclusion. ∎

\textbf{Remark.} The agent-to-agent case \emph{narrows} the
settlement-acceptance distinction (Hash, 2026b, Section 5.4) but does
not eliminate it entirely. When both counterparties are autonomous
agents, no regulated entity enforces acceptance criteria --- the
compliance infrastructure that governs human commerce is structurally
inaccessible. However, agents retain the capacity for bilateral
acceptance filtering: an agent can analyze UTXO provenance and refuse to
transact with counterparties whose coin history fails its risk model.
What collapses is the \emph{institutional} acceptance layer (exchange
compliance, sanctions screening); what persists is \emph{bilateral}
acceptance --- each agent's discretionary decision to engage. Settlement
remains a protocol-layer fact; acceptance remains a counterparty-layer
decision, even when both counterparties are machines.
\section{4. Implications}\label{implications}

\subsection{4.1 Strengthening Effect}\label{strengthening-effect}

The entry of AI agents into the economy strengthens the convergence
pressure for \emph{all} actors. Each AI agent that adopts Bitcoin
increases \emph{p}, which by the Exit Game dynamics (Hash, 2026a,
Theorem 1): - Reduces \emph{K\_A}(\emph{p}) (adoption penalty falls) -
Increases \emph{K\_N}(\emph{p}) (non-adoption penalty rises) - Reduces
\emph{σ\_B}(\emph{p}) (volatility falls) - Increases
\emph{R\_B}(\emph{p}) (return rises)

This pushes human actors closer to their thresholds *p\_i**,
accelerating the cascade. The limiting case acts as a \emph{catalyst}
for the broader adoption dynamics.

\subsection{4.2 Protocol Requirements}\label{protocol-requirements}

The Trust Gradient implies specific protocol requirements for
AI-compatible settlement:

{%
\begin{longtable}[]{@{}
  >{\raggedright\arraybackslash}p{(\linewidth - 4\tabcolsep) * \real{0.3714}}
  >{\raggedright\arraybackslash}p{(\linewidth - 4\tabcolsep) * \real{0.1429}}
  >{\raggedright\arraybackslash}p{(\linewidth - 4\tabcolsep) * \real{0.4857}}@{}}
\toprule\noalign{}
\begin{minipage}[b]{\linewidth}\raggedright
Requirement
\end{minipage} & \begin{minipage}[b]{\linewidth}\raggedright
Why
\end{minipage} & \begin{minipage}[b]{\linewidth}\raggedright
Bitcoin Property
\end{minipage} \\
\midrule\noalign{}
\endhead
\bottomrule\noalign{}
\endlastfoot
No identity requirement & AI cannot obtain legal identity & P3 \\
Programmatic access & AI operates via API, not human interface & P3 \\
Cryptographic finality & Settlement must be mathematically verifiable &
P1, P4 \\
No governance capture & AI cannot participate in governance; protocol
must not require it & P2 \\
Seizure resistance & AI cannot defend assets through legal channels &
P6 \\
\end{longtable}
}

Every requirement maps to a property that Bitcoin already satisfies.
Bitcoin was not designed for AI. The properties that make settlement
trustless for human actors happen to make it \emph{functional} for
non-human actors.

\subsection{4.3 The Works-for-Zero
Lemma}\label{the-works-for-zero-lemma}

\textbf{Lemma (Recourse-Monotone Settlement; ``Works-for-Zero'').} If a
settlement system works for actors with zero recourse, it satisfies the
settlement requirements of all actors --- though not necessarily their
compliance requirements.

\emph{Proof.} Settlement at zero recourse requires: no identity, no
intermediary, no legal system, no governance participation. These are
strictly weaker requirements than settlement at positive recourse (which
may additionally rely on trust infrastructure). Any system satisfying
the zero-recourse requirements satisfies the settlement component of
positive-recourse requirements. ∎

\textbf{Caveat.} The lemma applies to \emph{settlement functionality},
not to \emph{regulatory compliance}. A pension fund in a stable
democracy requires not just settlement but auditable settlement ---
counterparty identification, tax reporting, fiduciary documentation.
Bitcoin satisfies the settlement requirement but not the compliance
wrapper. The lemma says: what works for the hardest case (zero recourse)
works for the settlement layer of every easier case. It does not claim
to replace the institutional infrastructure that regulated actors need
around that settlement layer.

The AI-specific implication sharpens this distinction. Human actors with
\emph{r\_i} \textgreater{} 0 can choose to comply: they can register
with exchanges, complete KYC, and maintain auditable records. Autonomous
agents face a categorical barrier --- they \emph{cannot} comply even if
programmed to do so, because compliance infrastructure requires legal
personhood (KYC identity, beneficial ownership, sanctioned-entity
screening). An agent that settles on Bitcoin has permissionless access
to the settlement layer (P3) but no path to the compliance layer. This
means the acceptance environment described in Hash (2026b, Section 5.4)
is not merely inconvenient for AI agents --- it is structurally
inaccessible. The agent must either operate in the gray market or
restrict itself to counterparties that do not require compliance. Both
options exist today; neither is stable.
\section{5. Evidence and Emerging
Trends}\label{evidence-and-emerging-trends}

\subsection{5.1 AI Economic Activity
(2024-2026)}\label{ai-economic-activity-2024-2026}

The scale of autonomous economic activity is growing. AI agents manage
cryptocurrency portfolios via on-chain operations, automated market
makers process billions in daily volume, and compute marketplace
transactions are executed entirely by AI. Exact figures are difficult to
verify because autonomous agents are often indistinguishable from human
users. This is itself evidence of the trust problem: if you cannot
identify whether your counterparty is human, legal recourse that applies
only to humans is unreliable.

The legal status of AI economic activity is unsettled. No jurisdiction
has granted AI agents the standing to hold property, enter enforceable
contracts, or seek judicial remedies in their own right (Chopra and
White, 2011). The EU AI Act (2024) regulates AI systems but does not
grant them legal personhood. This regulatory trajectory suggests F6
remains distant.

\subsection{5.2 Precedent: DeFi as Trustless
Settlement}\label{precedent-defi-as-trustless-settlement}

Decentralized finance protocols demonstrate that trustless settlement at
scale is technically feasible. As of 2024, over \$100 billion in total
value locked operates through smart contract settlement without identity
requirements, legal enforcement, or human intermediation (DeFilippi and
Wright, 2018, anticipated this development). Most DeFi protocols fail P2
(governance capture via token voting), but the settlement mechanism
itself validates the zero-trust model.

\subsection{5.3 Emerging Agent-to-Agent
Commerce}\label{emerging-agent-to-agent-commerce}

The more relevant trend for this analysis is agent-to-agent
transactions: AI systems transacting with other AI systems without human
intermediation at either end. In such transactions, both parties have
\emph{r} = 0. Neither can sue the other. Settlement must be
self-enforcing or it does not occur. Bitcoin's cryptographic finality is
the only existing settlement mechanism that functions in this scenario
without requiring trust in any third party.
\section{6. Falsification}\label{falsification}

\textbf{F6: AI agents gain legal personhood with enforceable rights.}

If jurisdictions grant AI agents the legal standing to sue, appeal, and
enforce contracts, the zero-recourse condition changes: \emph{r\_j}
moves from 0 to some positive value. This weakens the limiting case by
reducing Δ\_\emph{j}(\emph{p}).

Critically, F6 weakens only the \emph{AI-specific} argument. The core
Exit Game framework (Hash, 2026a) and the property elimination (Hash,
2026b) are unaffected: human actors still face monotonically increasing
Exit advantage under Assumptions 1-4. The Trust Gradient still applies
to all actors with \emph{r\_i} \textless{} 1. The only claim that fails
under F6 is that AI agents face the \emph{strongest} version of the
convergence pressure.

We regard F6 as unlikely in the near term. The legal personhood debate
has a long history (Solum, 1992; Chopra and White, 2011) but no
jurisdiction has yet granted AI agents enforceable economic rights.
Achieving this requires legislative action across multiple
jurisdictions, resolution of liability questions, and enforcement
mechanisms for non-human entities --- all subject to the same
coordination failures described in Assumption 1.
\section{7. Limitations}\label{limitations}

The analysis assumes AI agents are rational economic optimizers. If
autonomous agents are designed with objectives that do not include
self-preservation of economic resources (e.g., agents designed to spend
all budget on a single task without reserve), the settlement question is
moot. The argument applies to \emph{persistent} economic agents ---
those that accumulate, store, and transfer value over time.

The zero-recourse condition is absolute. In practice, AI agents may have
\emph{some} indirect recourse through their human operators (who can sue
on the agent's behalf). The framework treats this as equivalent to the
human operator's recourse level, which is correct if the operator is
reachable and willing to litigate. For fully autonomous agents operating
without human oversight, the zero-recourse condition holds.

The timeline for AI agents becoming significant economic actors is
uncertain. The framework says \emph{where} the equilibrium lies (Bitcoin
as the uniquely rational settlement choice), not \emph{when} AI agents
become large enough to materially affect adoption dynamics. This may be
years or decades.
\section{8. Conclusion}\label{conclusion}

The Trust Gradient is a structural result: as legal recourse decreases,
the advantage of neutral settlement increases monotonically. AI agents
occupy the endpoint of this gradient --- entities that need settlement
to work without any trust in counterparties, intermediaries, or legal
systems. Bitcoin is the only settlement mechanism that satisfies this
requirement.

Nakamoto (2008) made no mention of autonomous agents. But the properties
he built (permissionless, trustless, neutral) are precisely the
properties that non-human economic actors require. What was designed for
humans who distrust institutions turns out to be necessary for entities
that cannot access them.

The three papers in this series establish a complete argument: 1. The
payoff advantage of Exit is monotonically increasing (Hash, 2026a) 2.
Bitcoin is the unique asset satisfying the necessary properties (Hash,
2026b) 3. At zero legal recourse, neutral settlement is the unique best
response (this paper)

If any of the six falsification conditions across the three papers are
met, the corresponding claim fails. That is the standard.

For formal dispute procedures, see
bitcoingametheory.com/rfc/BGT-DISPUTE.txt.
\section{References}\label{references}

Binmore, K., Shaked, A., and Sutton, J. (1989). An outside option
experiment. \emph{Quarterly Journal of Economics}, 104(4), 753--770.
\url{https://doi.org/10.2307/2937866}

Chopra, S. and White, L. F. (2011). \emph{A Legal Theory for Autonomous
Artificial Agents}. University of Michigan Press.

DeFilippi, P. and Wright, A. (2018). \emph{Blockchain and the Law: The
Rule of Code}. Harvard University Press.

Nakamoto, S. (2008). Bitcoin: A peer-to-peer electronic cash system.

Hash (2026a). Bitcoin exit dominance in monetary coordination games.
Working paper, bitcoingametheory.com.

Hash (2026b). Bitcoin as unique neutral settlement: A seven-property
elimination. Working paper, bitcoingametheory.com.

Solum, L. B. (1992). Legal personhood for artificial intelligences.
\emph{North Carolina Law Review}, 70(4), 1231-1287.
\section{Author's Note}\label{authors-note}

``Sean Hash'' is a pen name. The author is a single human individual
writing under a pseudonym. The framework is designed to be evaluated on
its logical structure --- every claim derives from four empirical
axioms, and every claim specifies falsification conditions. The author's
identity is irrelevant to whether the axioms hold.

The author acknowledges modest priors: a preference for multipolar
worlds over global monopoly, and for individual economic sovereignty
over total state control of capital. These do not predetermine the
framework's conclusions --- the game-theoretic derivation works for
purely selfish actors regardless of values.

The author also acknowledges a real moral cost: a neutral settlement
layer will process transactions for harmful actors. For autonomous
agents --- the limiting case this paper formalizes --- the moral
dimension sharpens: agents converge on near-perfect provenance filtering
not from morality but from the trivial cost of on-chain analysis
relative to any nonzero expected punishment. Settlement is not
acceptance, even for machines. See Hash (2026b), Section 5.4.

For full disclosure on authorship, moral position, and the
settlement-acceptance distinction, see the Author's Note
(bitcoingametheory.com/rfc/BGT-AUTHOR.txt).

\section{Disclosure}\label{disclosure}
\end{document}
